\input{preamble}

\title{Section 1.8 --- Absolute Value Equations}

\begin{document}
\maketitle

%\section{Evaluating Absolute Values}
%Evaluate each expression.
%\begin{smenum}
%\mitemxx
%{$\Abs{2-3}+\dfrac{\Abs{5}}{\Abs{-4}}$}
%{$-\Abs{-5}\cdot\Abs{7-1}$}
%\mitemxx
%{$2\cdot-\Abs{6+3}-\Abs{-2\cdot 6}$}
%{$-\frac{1}{3}\Abs{2-8}+4$}
%\end{smenum}
%
%%%%%%%%%%%%%%%%%%%%%%%%%%%%%%%%%%%%%%%%%%%%%%%%%%%%%%%%%
%\section{Identifying Absolute Value Functions}
%Determine whether each of the following is an absolute value function.
%\begin{smenum}
%\mitemxxxx
%{$f(x)=\Abs{x-1}+3$}
%{$g(x)=-(2x+4)+6$}
%{$h(x)=\Abs{x^2-3x+2}$}
%{$k(x)=\frac{1}{2}\Abs{4-x}+1$}
%\mitemxxxx
%{$r(x)=\dfrac{3}{\Abs{4x-1}}$}
%{$s(x)=\sqrt{\Abs{2x-1}}$}
%{$t(x)=-\Abs{2x-3}-\frac{3}{4}$}
%{$u(x)=\Abs{\frac{1}{3}x-5}-\dfrac{4}{x^2}$}
%\end{smenum}
%
%%%%%%%%%%%%%%%%%%%%%%%%%%%%%%%%%%%%%%%%%%%%%%%%%%%%%%%%%
%\section{Graphing Absolute Value Functions}
%Graph each absolute value function. Then, looking at the graph, state the domain and range of each function.
%\begin{smenum}
%\mitemxx
%{$f(x)=\frac{2}{3}\Abs{x}$\\
%\begin{ilpic}[x=0.2in,y=0.2in]
%%\useasboundingbox(-5.9,-5.9) rectangle (5.9,5.9);
%\DrawGridSquare
%\end{ilpic}}
%{$g(x)=2\Abs{x-4}-2$\\
%\begin{ilpic}[x=0.2in,y=0.2in]
%%\useasboundingbox(-5.9,-5.9) rectangle (5.9,5.9);
%\DrawGridSquare
%\end{ilpic}}
%\mitemxx
%{$h(x)=-\frac{3}{4}\Abs{x}+1$\\
%\begin{ilpic}[x=0.2in,y=0.2in]
%%\useasboundingbox(-5.9,-5.9) rectangle (5.9,5.9);
%\DrawGridSquare
%\end{ilpic}}
%{$k(x)=-\frac{3}{2}\Abs{x+1}-1$\\
%\begin{ilpic}[x=0.2in,y=0.2in]
%%\useasboundingbox(-5.9,-5.9) rectangle (5.9,5.9);
%\DrawGridSquare
%\end{ilpic}}
%\end{smenum}
%
%%%%%%%%%%%%%%%%%%%%%%%%%%%%%%%%%%%%%%%%%%%%%%%%%%%%%%%%%
%\section{Range and Domain of an absolute value function}
%Find the domain and range of each function.
%\begin{smenum}
%\mitemxx
%{$f(x)=\Abs{x-2}$}
%{$g(x)=2\Abs{x}+4$}
%\mitemxx
%{$h(x)=-3\Abs{x-1}+3$}
%{$k(x)=\frac{1}{3}\Abs{2x+1}-4$}
%\end{smenum}

%%%%%%%%%%%%%%%%%%%%%%%%%%%%%%%%%%%%%%%%%%%%%%%%%%%%%%%%
\section{Zero, One, or Two Solutions}
Determine whether each absolute value equation has zero, one, or two solutions.  If the equation has one or two solutions, find them.
\begin{smenum}
\mitemxxxx
{$\Abs{x}=4$}
{$\Abs{2x}=-2$}
{$\Abs{3x+18}=0$}
{$\Abs{x-6}=17$}
\mitemxxxx
{$\Abs{4x-6}=-2$}
{$\Abs{\frac{2}{3}x+4}=5$}
{$\Abs{\frac{4}{3}x-\frac{1}{3}}=0$}
{$\Abs{\frac{2}{5}x+\frac{5}{3}}=-\frac{2}{13}$}
\end{smenum}

%%%%%%%%%%%%%%%%%%%%%%%%%%%%%%%%%%%%%%%%%%%%%%%%%%%%%%%%
\section{Solve Absoulte Value Equations}
Find any possible solutions to each equation.
\begin{smenum}
\mitemxx
{$\Abs{3r}+7=10$}
{$-4\Abs{x+3}-1=-9$}
\mitemxx
{$\Abs{v-2}+4=2$}
{$3=1-\Abs{4u-8}$}
\mitemxx
{$3-2\Abs{x-5}=-7$}
{$2\Abs{4y-2(6-y)}=36$}
\mitemxx
{$\Abs{\frac{u}{3}-1}=4$}
{$-3\Abs{p+\frac{2}{3}}-1=-10$}
\mitemxx
{$5-3\Abs{2x+3}=5$}
{$\frac{2}{3}\Abs{7-\frac{1}{3}x}=6$}
\end{smenum}

\end{document}